\chapter{From Transformer to GPT}
\label{ch:from-transformer-to-gpt}

\abstract{This chapter specializes the \transformer into a causal decoder. It covers next-token loss, GPT-style blocks, sampling, repetition control, loss curves, minimal training loops, and efficient attention as used in practical GPT implementations.}

\section{Chapter Spine}
GPT is not merely a \transformer with a different name; it is a causal training and generation interface. The chapter moves from the decoder block to cross-entropy training, then from logits to sampling behavior.

\section{Core Sections}
\subsection{Causal Language Modeling}
Teacher forcing, shifted labels, cross-entropy, perplexity, and batch packing.

\subsection{Minimal GPT Implementation}
Configuration, block stack, tied embeddings, optimizer groups, and checkpointing.

\subsection{Decoding}
Greedy decoding, temperature, top-$k$, nucleus sampling, repetition penalties, and stop conditions.

\subsection{Attention Efficiency}
FlashAttention, memory bandwidth, IO-awareness, and why exact attention can be made faster without changing model semantics \cite{dao2022flashattention}.

\section{Planned Exercises}
Train a toy character-level model, reproduce a loss curve, and compare sampling controls on the same prompt.
